% Generated draft: requires mathematical and source review.
\documentclass[11pt]{article}
\usepackage[margin=25mm]{geometry}
\usepackage{fontspec}
\setmainfont{Latin Modern Roman}
\usepackage{amsmath,amssymb,mathtools,unicode-math}
\setmathfont{Latin Modern Math}
\usepackage{enumitem,needspace,xurl,hyperref}
\setlength{\parindent}{0pt}
\setlength{\parskip}{5pt}
\setlength{\emergencystretch}{3em}
\newcommand{\N}{\mathbb{N}}
\newcommand{\Z}{\mathbb{Z}}
\newcommand{\Q}{\mathbb{Q}}
\newcommand{\R}{\mathbb{R}}
\newcommand{\C}{\mathbb{C}}
\newcommand{\F}{\mathbb{F}}
\newcommand{\A}{\mathbb{A}}
\newcommand{\PP}{\mathbb{P}}
\newcommand{\E}{\mathbb{E}}
\newcommand{\eps}{\varepsilon}
\newcommand{\dd}{\,\mathrm{d}}
\begin{document}
\section*{UnsolvedMath ID: 11000267}
\textbf{Source problem code:} catalogue entry
\subsection{Definitions and mathematical statement}
For \(n\geq2\), Culler--Vogtmann Outer space in its usual projectivized model consists of finite connected rank-\(n\) core metric graphs of total edge length \(1\), with positive edge lengths and a homotopy-equivalence marking by the rank-\(n\) rose, modulo marking-preserving isometry. It carries the \(\operatorname{Out}(F_n)\)-action. The problem asks to develop a metric theory of Outer space; the source does not prescribe its variant or a symmetric/asymmetric metric.

An iwip, or fully irreducible outer automorphism, has no positive power preserving the conjugacy class of a proper nontrivial free factor. The accompanying discussion seeks an axis theory for such elements, motivated by axial semisimple integral-matrix actions on the source's symmetric space \(X\) and unique pseudo-Anosov axes in Teichm\"uller space. It reports that iwips and pseudo-Anosovs have two distinguished natural-boundary fixed points with north--south dynamics for their cyclic actions [62, 54]. The exact compactifications and metric hypotheses are not supplied.

The source says initial Handel--Mosher work [43] rules out expecting a unique metric axis for iwips, but later work [42] directly constructs an axis bundle for each iwip, with comparison [63]. Thus it retains the broader metric-theory problem despite that progress, rather than asking only for an axis already supplied by the later construction. The claimed matrix analogy, including \(X=\mathbb H^2\) in rank two, is contextual and does not specify a chosen Outer-space metric or rank-one property in higher rank.

The extract ends by announcing a modified question, without printing it. Therefore the desired specific remaining metric properties, weaker-axis axioms or quantitative bounds are not recoverable. No conjectured relative hyperbolicity or unique-axis theorem is imposed.
\subsection{Short English statement}
Develop an Outer-space metric and axis theory compatible with fully irreducible dynamics, retaining Handel--Mosher axis-bundle progress and the missing final formulation.
\subsection{Sources}
\begin{itemize}
\item \emph{Problems on Mapping Class Groups and Related Topics}, edited by Farb, Question 1.2 and surrounding metric-theory discussion, supplied PDF page 328. \url{https://www.math.uchicago.edu/~farb/papers/mcgbook.pdf}. Recorded access 2026-07-29; extraction pdf-font-structure.
\item Handel and Mosher [43] and [42], with [63] for axis-bundle comparison; [62], [54] for the boundary-dynamics discussion. Complete bibliography entries are absent.
\item Local UnsolvedMath registry, \texttt{data/unsolvedmath/problems.json}, permanent ID \(11000267\), exact source code \texttt{catalogue entry}. Its OPEN-TRIAGE summary mentions Handel--Mosher and relative-hyperbolicity work without a precise metric, model or reference.
\end{itemize}
\end{document}
